\documentclass[10pt,a4paper]{amsart}
\usepackage[margin=19mm]{geometry}
\usepackage{amsmath,amssymb,mathtools}
\usepackage[scr=rsfs,cal=euler]{mathalfa}
\usepackage{xfrac}
\usepackage[unicode,hidelinks,bookmarks=false]{hyperref}
\newcommand{\ie}{i.e.\ }
\newcommand{\cf}{cf.\ }
\newcommand{\resp}{respectively\ }
\newcommand{\Subsection}{\subsection}
\providecommand{\nobreakdash}{\nobreak\hbox{-}}
\setlength{\emergencystretch}{2em}
\multlinegap0pt
\usepackage{amssymb}
\usepackage{bm}
\usepackage{xcolor}
\usepackage{float}
\usepackage[shortlabels]{enumitem}
\newtheorem{proposition}{Proposition}
\title{A Minimal Four-Thruster System for Comet-Based Interstellar Navigation}
\author{Bo Pieter Johannes Andr\'ee}
\date{Source: arXiv:2603.16981v3 (CC BY 4.0)}
\begin{document}
\maketitle

\noindent\emph{Source and licence.} A Minimal Four-Thruster System for Comet-Based Interstellar Navigation. Version arXiv:2603.16981v3 (CC BY 4.0), distributed by arXiv under the Creative Commons Attribution 4.0 International licence, http://creativecommons.org/licenses/by/4.0/, as recorded in the arXiv licence metadata for this exact version and shown on the abstract page https://arxiv.org/abs/2603.16981v3. Full licence notice: \texttt{licenses/arxiv-2603.16981-CC-BY-4.0.txt}.\par\smallskip

\noindent\textit{Editorial scope.} Counted unit: the article's proposition prop:thrust\_neutral\_unique (source0.tex lines 1124--1142). The article's complete original proof of it is delivered with it (source0.tex lines 1144--1206, one \texttt{proof} environment of 63 lines). No mathematics is written, repaired or re-derived: the statement and the proof are the article's own lines, in the article's own notation, and no retained line is substituted or re-flowed.  Retained uncounted as the source-local context the proof needs: lines 1047--1058.  Nothing was omitted from any retained span, so the package carries no omission record.  No retained line contains a citation, so the wrapper carries no bibliography and no bibliography entry.  No retained line contains a figure environment, an image-inclusion command or drawing code, so no image is delivered and the package declares no asset dependency.  Editorial formatting only, in addition: an AMS \texttt{amsart} preamble that restates the article's own declarations and macros used by the retained lines, namely the environments \texttt{proposition} on separate counters, together with the standard packages those lines need.  The retained environments print in this document as Proposition 1 in order of appearance, whereas the article prints the same items with the numbering of its own full document.  The complete native source of the article as distributed in the arXiv e-print for this exact version is bundled unchanged as \texttt{sources/196588/source0.tex}.
\begin{proposition}[Differential steering decomposition]
\label{prop:differential_decomposition}
Let $u_i = u_0 + \delta_i$ with baseline $u_0 \ge 0$ and differentials $\delta_i \in \mathbb{R}$ 
satisfying $u_i \ge 0$ for all $i$. Then the net force and total thrust decompose as
\begin{equation}
\mathbf F = \sum_{i=1}^{3} \delta_i \mathbf d_i,
\qquad
S = 3u_0 + \sum_{i=1}^{3} \delta_i.
\end{equation}
In particular, the steering force depends only on the differentials $\delta_i$, and the change 
in total thrust relative to the symmetric baseline is $\Delta S = \sum_{i=1}^{3} \delta_i$.
\end{proposition}

\begin{proposition}[Thrust-neutral steering: existence and uniqueness]
\label{prop:thrust_neutral_unique}
For any desired in-plane force $\mathbf F = (F_x, F_y)^\top$, there exists a \emph{unique} 
differential vector $\bm\delta = (\delta_1, \delta_2, \delta_3)^\top$ satisfying the 
\emph{thrust-neutral constraint}
\begin{equation}
\delta_1 + \delta_2 + \delta_3 = 0
\label{eq:thrust_neutral}
\end{equation}
and generating $\mathbf F$ via $\sum_{i=1}^{3} \delta_i \mathbf d_i = \mathbf F$. 
The unique solution is
\begin{align}
\delta_1 &= \tfrac{2}{3} F_x, \label{eq:delta1} \\
\delta_2 &= -\tfrac{1}{3} F_x + \tfrac{1}{\sqrt{3}} F_y, \label{eq:delta2} \\
\delta_3 &= -\tfrac{1}{3} F_x - \tfrac{1}{\sqrt{3}} F_y. \label{eq:delta3}
\end{align}
Consequently, thrust-neutral steering changes the net force without changing total thrust: 
$S$ remains equal to $3u_0$.
\end{proposition}

\begin{proof}
Write the force and neutrality conditions as a $3 \times 3$ linear system in 
$(\delta_1, \delta_2, \delta_3)$:
\begin{align}
F_x &= \delta_1 - \tfrac{1}{2}\delta_2 - \tfrac{1}{2}\delta_3, \label{eq:sys_x} \\
F_y &= \tfrac{\sqrt{3}}{2}\delta_2 - \tfrac{\sqrt{3}}{2}\delta_3, \label{eq:sys_y} \\
0   &= \delta_1 + \delta_2 + \delta_3. \label{eq:sys_neut}
\end{align}

\emph{Step 1: Solve for $\delta_1$.}
Rearrange~\eqref{eq:sys_neut} to obtain $\delta_2 + \delta_3 = -\delta_1$. 
Substitute into~\eqref{eq:sys_x}:
\[
\begin{aligned}
F_x &= \delta_1 - \tfrac{1}{2}(\delta_2 + \delta_3) \\
    &= \delta_1 - \tfrac{1}{2}(-\delta_1) \\
    &= \tfrac{3}{2}\delta_1.
\end{aligned}
\]
Solving yields $\delta_1 = \tfrac{2}{3} F_x$.

\emph{Step 2: Solve for $\delta_2$ and $\delta_3$.}
From~\eqref{eq:sys_y}, divide by $\tfrac{\sqrt{3}}{2}$ to get
\[
\delta_2 - \delta_3 = \tfrac{2}{\sqrt{3}} F_y.
\]
We now have two equations in $\delta_2, \delta_3$:
\[
\begin{aligned}
\delta_2 + \delta_3 &= -\tfrac{2}{3} F_x, \\
\delta_2 - \delta_3 &= \tfrac{2}{\sqrt{3}} F_y.
\end{aligned}
\]
Add these equations and divide by 2:
\[
\begin{aligned}
\delta_2 &= \tfrac{1}{2}\Bigl(-\tfrac{2}{3} F_x + \tfrac{2}{\sqrt{3}} F_y\Bigr) \\
         &= -\tfrac{1}{3} F_x + \tfrac{1}{\sqrt{3}} F_y.
\end{aligned}
\]
Subtract and divide by 2:
\[
\begin{aligned}
\delta_3 &= \tfrac{1}{2}\Bigl(-\tfrac{2}{3} F_x - \tfrac{2}{\sqrt{3}} F_y\Bigr) \\
         &= -\tfrac{1}{3} F_x - \tfrac{1}{\sqrt{3}} F_y.
\end{aligned}
\]

\emph{Uniqueness.}
The coefficient matrix of system~\eqref{eq:sys_x}--\eqref{eq:sys_neut} is
\[
A = \begin{pmatrix}
1 & -\tfrac{1}{2} & -\tfrac{1}{2} \\[2pt]
0 & \tfrac{\sqrt{3}}{2} & -\tfrac{\sqrt{3}}{2} \\[2pt]
1 & 1 & 1
\end{pmatrix}.
\]
A direct calculation gives $\det(A) = \tfrac{3\sqrt{3}}{2} \neq 0$, so the system has a unique 
solution for any $(F_x, F_y)$.

Finally, since $\sum_{i=1}^{3} \delta_i = 0$, Proposition~\ref{prop:differential_decomposition} 
implies $S = 3u_0 + 0 = 3u_0$.
\end{proof}

\end{document}
